%%
%% This is file `mathfont_heading.tex',
%% generated with the docstrip utility.
%%
%% The original source files were:
%%
%% mathfont_code.dtx  (with options: `heading')
%% 
%% This file is from version 2.0 of the free and open-source
%% LaTeX package "mathfont," to be used with the XeTeX or
%% LuaTeX engines. (As of version 2.0, LuaTeX is recommended.)
%% 
%% Copyright 2018-2021 by Conrad Kosowsky
%% 
%% This file may be distributed and modified under the terms
%% of the LaTeX Public Project License, version 1.3c or any
%% later version. The most recent version of this license is
%% available online at
%% 
%%           https://www.latex-project.org/lppl/.
%% 
%% This Work has the LPPL status "maintained," and the current
%% maintainer is the package author, Conrad Kosowsky. He can
%% be reached at kosowsky.latex@gmail.com.
%% 
%% THIS SOFTWARE IS PROVIDED "AS IS" WITHOUT ANY EXPRESS
%% OR IMPLIED WARRANTY, INCLUDING THE IMPLIED WARRANTIES
%% OF MERCHANTABILITY OR FITNESS FOR A PARTICULAR PURPOSE.
%% IF THE SOFTWARE IS DEFECTIVE, YOU AGREE TO ASSUME THE
%% COST FOR ANY REPAIR OR CORRECTION.
%% 
%% BY USING OR DISTRIBUTING THIS SOFTWARE, YOU AGREE
%% TO RELEASE THE PACKAGE AUTHOR FROM ANY LIABILITY FOR
%% DAMAGES ARISING OUT OF YOUR USE OF THE SOFTWARE. YOU
%% AGREE TO DO SO TO THE MAXIMUM EXTENT ALLOWED UNDER
%% APPLICABLE LAW AND EVEN IF THE PACKAGE AUTHOR HAS
%% BEEN ADVISED OF THE POSSIBILITY OF SUCH DAMAGES.
%% 
%% See also the "No Warranty" section of the LaTeX Project
%% Public License. In releasing the package author from
%% liability, you also release from liability any third
%% parties who distribute the software under the terms
%% of the LaTeX Project Public License. Derivative works
%% based on this package may come with their own license or
%% terms of use, and the package author is not responsible
%% for any third-party software.
%% 
%% For more information, see the mathfont_code.dtx.
%% 
%% Happy TeXing!
%% 

\def\packagedate{December 2021}
\def\packageversion{2.0}


\let\@@section\section
\def\section{\@ifstar\star@sect\no@star@sect}
\def\star@sect#1{\@@section*{#1}\section@name{#1}}
\def\no@star@sect#1{\@@section{#1}\section@name{#1}}
\def\section@name#1{\expandafter
  \def\csname section@\thesection\endcsname{#1}}
\def\sectionname{\csname section@\thesection\endcsname}

\def\@oddhead{\ifnum\count0>1\relax
  \rlap{\textit{\sectionname}}\hfil
  \hbox to 0pt{\hss\documentname\hss}\hfil
  \llap{\the\count0}\fi}
\def\@evenhead{\ifnum\count0>1\relax
  \rlap{\the\count0}\hfil
  \hbox to 0pt{\hss\documentname\hss}\hfil
  \llap{\textit{\sectionname}}\fi}
\def\@oddfoot{\hfil\ifnum\count0=1\relax1\fi\hfil}
\let\@evenfoot\@empty

\pretolerance=20
\hyphenpenalty=10
\exhyphenpenalty=5
\brokenpenalty=0
\clubpenalty=5
\widowpenalty=5
\finalhyphendemerits=300
\doublehyphendemerits=500

\protected\def\XeTeX{X\kern-0.1em
  \raise-0.5ex\hbox{\rotatebox[origin=c]{180}{E}}\kern-0.15em
  \TeX}
\protected\def\XeLaTeX{X\kern-0.1em
  \raise-0.5ex\hbox{\rotatebox[origin=c]{180}{E}}\kern-0.13em
  \LaTeX}
\bgroup
  \count@\catcode`\|
  \catcode`\|=12\relax
  \gdef\indexpage#1{\index{#1|usage}}
\egroup
\protected\def\usage#1{\textit{#1}}
\bgroup
\catcode`\_=12
  \protected\gdef\fontspeccommand{\texttt{\string\fontspec_set_family:Nnn}}
  \protected\gdef\fontspecbool{\texttt{\string\g__fontspec_math_bool}}
\egroup
\def\topfraction{1}
\def\bottomfraction{1}
\newenvironment{code}
  {\strut\vadjust\bgroup\medskip\parindent\z@\leftskip2em\relax
    \noindent\strut\ignorespaces}
  {\strut\par\medskip\egroup\hfill\break\strut\ignorespacesafterend}
\def\vrb#1{\expandafter\texttt\expandafter{\string#1}}
\parskip=0pt

\def\charexample#1{\hbox to \hsize{\hbox to 0.4in{$#1$\hfil}\vrb#1\hfil}}
\def\accentexample#1{\hbox to \hsize{\hbox to 0.4in{$#1 a$\hfil}\vrb#1\hfil}}
\def\delimexample#1{\hbox to \hsize{%
  \hbox to 0.8in{$#1\big#1\Big#1\bigg#1\Bigg#1$\hfil}\vrb#1\hfil}}
\def\luadelimexample#1{\hbox{\vbox{\hbox to 0.8in{%
    $#1\big#1\Big#1\bigg#1\Bigg#1$\hfil}\hrule width 0pt height 0pt\relax}%
  \hbox to \dimexpr\hsize-0.8in{\vbox{\vss
    \hbox{\vrb#1}
    \hbox{\quad(Lua\TeX only)}
    \vss}\hfil}}}
\def\operatorexample#1{\hbox to \hsize{%
  \hbox to 0.4in{$#1$\hfil}%
  \hbox to 0.4in{$\displaystyle#1$\hfil}\vrb#1\hfil}\smallskip}
\def\blockheader#1#2{%
  \bigskip\medskip
  \penalty-9000
  \hbox{\fbox{\hbox to \dimexpr\hsize-2\fboxrule-2\fboxsep\relax{%
    \hfil#2 Characters (Keyword \texttt{#1})\strut\hfil}}}
  \medskip}
\def\upperalphabet{ABCDEFGHIJKLMNOPQRSTUVWXY}
\def\loweralphabet{abcdefghijklmnopqrstuvwxy}
\def\digits{0123456789}
\def\printchars#1{%
  \expandafter\@tfor\expandafter\letter\expandafter:\expandafter=#1\do{%
  \rlap{$\@tempstyle{\letter}$}\hfill}}
\def\letterlikechars#1{\smallskip\let\@tempstyle#1
  \noindent\printchars\upperalphabet\hbox to 0.6em{$\@tempstyle{Z}$\hss}\par
  \noindent\printchars\loweralphabet\hbox to 0.6em{$\@tempstyle{z}$\hss}\par}

{\large\parindent=0pt\leftskip=0pt plus 1 fil\rightskip=0pt plus 1fil\parfillskip=0pt
{\strut\Large Package \textsf{mathfont} v.\ \packageversion\ \documentname\let\thefootnote\relax\footnote{Acknowledgements: Thanks to Lyric Bingham for her work checking my unicode hex values. Thanks to Herbert Voss and Andreas Zidak for pointing out bugs in previous versions of \textsf{mathfont}. Thanks to Jean-Fran\c cois Burnol for pointing out an error in the documentation in reference to their \textsf{mathastext} package.}\global\advance\c@footnote\m@ne}\par
{\strut Conrad Kosowsky}\par
{\strut\packagedate}\par
{\strut\ttfamily kosowsky.latex@gmail.com}\par}

\bigskip

\ifnum\showabstract=1\relax

\hrule height \p@\hbox{\vrule width \p@\kern-\p@\relax\vbox{\medskip
{\leftskip=1.4in\rightskip=1.4in
\noindent\strut For easy, off-the-shelf use, type the following in your preamble and compile with \XeLaTeX\ or Lua\LaTeX:\par
\medskip
\vbox{\noindent\hfil{|\usepackage[|\meta{font name}|]{mathfont}|}\hfil}
\medskip
\noindent As of version 2.0, using Lua\LaTeX\ is recommended.\par}
\medskip}\kern-\p@\vrule width \p@}\hrule height \p@

\bigskip

{\small

\centerline{\bfseries Overview\strut}
\smallskip

\leftskip=0.5in\relax
\rightskip=0.5in\relax
\noindent The \textsf{mathfont} package adapts unicode text fonts for math mode. The package allows the user to specify a default unicode font for different classes of math symbols, and it provides tools to change the font locally for math alphabet characters. When typesetting with Lua\TeX, \textsf{mathfont} adds resizable delimiters, big operators, and a MathConstants table to text fonts.\par}

\bigskip\bigskip\nointerlineskip
\centerline{\vrule height 0.5pt width 2.5in}\bigskip\medskip
\fi
\endinput
%%
%% End of file `mathfont_heading.tex'.
